%% LyX 2.4.4 created this file.  For more info, see https://www.lyx.org/.
%% Do not edit unless you really know what you are doing.
\documentclass[english]{article}
\usepackage[T1]{fontenc}
\usepackage[latin9]{inputenc}
\usepackage{enumitem}
\usepackage{amsmath}
\usepackage{amsthm}
\usepackage{amssymb}
\usepackage{cancel}
\usepackage{geometry}
\geometry{verbose,tmargin=1in,bmargin=1in,lmargin=1in,rmargin=1in}

\makeatletter
%%%%%%%%%%%%%%%%%%%%%%%%%%%%%% Textclass specific LaTeX commands.
\numberwithin{equation}{section}
\numberwithin{figure}{section}
\newlength{\lyxlabelwidth}      % auxiliary length 

%%%%%%%%%%%%%%%%%%%%%%%%%%%%%% User specified LaTeX commands.
\usepackage{fancyhdr}
\usepackage{lastpage}
\pagestyle{fancy}
\usepackage{enumitem}
\usepackage{ifthen}
%sets math fonts to be more readable
\everymath=\expandafter{\the\everymath\displaystyle}
%\DeclareMathSizes{10}{10}{10}{10}
\usepackage{multicol}

\usepackage{tikz}
\usetikzlibrary{plotmarks}
\usepackage{pgfplots}
\pgfplotsset{compat=newest}

\usepackage{calc}
\usepackage[nomessages]{fp}

\usepackage{advdate}    % Advancing/saving dates
\usepackage{datetime}   % Dates formatting
\usepackage{datenumber} % Counters for dates
% Added by lyx2lyx
\setlength{\parskip}{\smallskipamount}
\setlength{\parindent}{0pt}

\makeatother

\usepackage{babel}
\begin{document}
\renewcommand{\author}{Dr.\,James Ashe}

\newcommand{\classNum}{131}

\newcommand{\classSec}{A}

\newcommand{\nameType}{}

\newcommand{\examType}{Exam 1 Review}

\renewcommand{\date}{\formatdate{11}{9}{2013}}

%%First page's header%%%%%%%%%%%%%%%%%%%%%%
\fancypagestyle{firststyle} %%header settings for first pagr
{
	\fancyhead[L]{MTH \classNum{}(\classSec{}) \author{}\\
	\textbf{\examType{}} \date{}}
	\fancyhead[C]{\textbf{\nameType}}
	\setlength{\headsep}{14pt}
}
%%%%%%%%%%%%%%%%%%%%%%%%%%%%%%%%%%%%%%%%%%

%%Next page's header%%%%%%%%%%%%%%%%%%%%%%%
\setlist{nolistsep}
\lhead{\classNum{}(\classSec{})} 
\chead{\textbf{\nameType}} 
\cfoot{\thepage{}~of~\pageref{LastPage}} 
\setlength{\headsep}{5pt} 
%%%%%%%%%%%%%%%%%%%%%%%%%%%%%%%%%%%%%%%%%%%

%%Various settings; see the preamble for other settings%%%%%
%%\setenumerate[0]{leftmargin=12pt,itemindent=0pt,labelwidth=0pt}
\renewcommand{\labelenumi}{\textbf{\arabic{enumi}.}}
\renewcommand{\labelenumii}{\textbf{\alph{enumii}.}}
\thispagestyle{firststyle}
%%%%%%%%%%%%%%%%%%%%%%%%%%%%%%%%%%%%%%%%%%
\newcommand{\xmax}{10}
\newcommand{\xmin}{-10}
\newcommand{\xstep}{0} %must be xmin+n
\newcommand{\ymax}{10}
\newcommand{\ymin}{-10}
\newcommand{\ystep}{0} %must be ymin+n

\newcommand{\Dspace}{\vspace{1cm}}\textbf{Plotting point on the graph and reading points from the graph}
\begin{enumerate}
\item Given a point $(x,y)$ count $x$ \emph{horizontally }in the direction
of the sign and count $y$ \emph{vertically }in the direction of the
sign.
\end{enumerate}
\textbf{Find the $x$ and $y$ intercepts from a graph. }(See the
midterm review for examples)
\begin{enumerate}
\item Find all the times the graph intersects the $x-$axis. These are your
$x-$intercepts.
\item Find all the times the graph intersects the $y-$axis. These are your
$y-$intercepts.
\item [\textbf{ex.}]\vspace{0cm}\renewcommand{\xmax}{10}
\renewcommand{\xmin}{-10}
\renewcommand{\xstep}{0} %must be xmin+n
\renewcommand{\ymax}{10}
\renewcommand{\ymin}{-10}
\renewcommand{\ystep}{0} %must be ymin+n
%%%%%%%
\begin{tikzpicture} 
\begin{axis}[height=5.5cm, 
	width=5.5cm, 
	axis lines=middle,
	minor x tick num={4},
	minor y tick num={4},
	grid=both,
	%xtick={0,50,...,250},
	%ytick={0,10,20,...,40},
	%axis line style={-latex}, 
	xmin=\xmin,xmax=\xmax, 
	ymin=\ymin,ymax=\ymax,
	%restrict y to domain=-1:10,
	%no markers, 
	xlabel={$x$},ylabel={$y$}, 
	%every axis x label/.append style={anchor=north}, 
	%every axis y label/.append style={anchor=east}
	]
	
	\addplot[color=blue,solid,mark=*] coordinates {(-3,-5)} node[pos=.65,above] {$(-3,5)$};
	\addplot[domain=\xmin:\xmax,thick,samples=100,draw=red,<->]{-1/2*x+3} node[pos=.65,above] {$f(x)$};
	
\end{axis}
\end{tikzpicture}
\end{enumerate}

\subsubsection*{Solve an equation. ex. $3(x+8)-2=-(x+3)-4x$.}
\begin{enumerate}
\item Simplify, distribute, and then simplfiy. Note that $-(x+3)=-1(x+3)$.
\begin{eqnarray*}
3x+24-2 & = & -x-3\\
3x+22 & = & -x-3
\end{eqnarray*}
\item Combine like terms and solve, i.e., put ``$x's$'' on one side and
numbers on the other..
\begin{eqnarray*}
4x & = & -25\\
x & = & \frac{-25}{4}
\end{eqnarray*}
\end{enumerate}

\subsubsection*{Solve an equation with fractions. ex. $\frac{2x}{6}+2=\frac{3x}{9}$}
\begin{enumerate}
\item Multiply by the Lowest Common Denominator (LCD). In this case, $18$.
Note BOTH terms on the left got multiplied by $18$.
\begin{eqnarray*}
18\left(\frac{2x}{6}+2\right) & = & \left(\frac{x}{9}\right)18\\
\frac{36x}{6}+36 & = & \frac{18x}{9}\\
6x+36 & = & 2x
\end{eqnarray*}
\item Now just solve.
\begin{eqnarray*}
4x & = & -36\\
x & = & -\frac{36}{4}=-9
\end{eqnarray*}
\end{enumerate}

\subsubsection*{Write the value(s) that make the denominator(s) zero. Then solve
an equation. ex. $\frac{5}{x-2}+\frac{2}{x+7}=\frac{10}{(x+7)(x-2)}$}
\begin{enumerate}
\item First note that $x\neq2,-7$.
\item Multiply by the LCD. In this case $(x-2)(x+7)$.
\begin{eqnarray*}
(x+7)(x-2)\left(\frac{5}{x-2}+\frac{2}{x+7}\right) & = & \left(\frac{10}{(x+7)(x-2)}\right)(x+7)(x-2)\\
5(x+7)+2(x-2) & = & 10
\end{eqnarray*}
\item Now solve. 
\begin{eqnarray*}
5(x+7)+2(x-2) & = & 10\\
7x+31 & = & 10\\
x & = & -3
\end{eqnarray*}
\end{enumerate}

\subsubsection*{Solve a formula for a specified variable. ex. $P=2L+2W$ for $L$.}
\begin{enumerate}
\item Just solve like you do for an $x$, and treat the other letters like
numbers. 
\begin{eqnarray*}
P & = & 2L+2W\\
P-2W & = & 2L\\
\frac{P-2W}{2} & = & L
\end{eqnarray*}
\item Note: The $2$ divides BOTH factors, and the you DO NOT cancel anything
in $2W$.
\end{enumerate}

\paragraph{Solve using the square root property. ex. $(x-4)^{2}+1=2$.}
\begin{enumerate}
\item Get the squared part by itself:
\[
(x-4)^{2}=1
\]
\item Use a square root; Don't forget the $\pm$. 
\begin{eqnarray*}
\sqrt{(x-4)^{2}} & = & \pm\sqrt{1}\\
x-4 & = & \pm1
\end{eqnarray*}
\item Solve.
\begin{eqnarray*}
x & = & \pm1+4\\
x=5 & \mbox{ and } & x=3
\end{eqnarray*}
\end{enumerate}

\subsubsection*{Solve by factoring. ex. $3x^{2}=-3x+6$}
\begin{enumerate}
\item Set the equation equal to $0$. 
\[
3x^{2}+x-2=0
\]
\item Factor any common multiples.
\[
3(x^{2}+x-2)=0
\]
\item Factor and solve.
\begin{eqnarray*}
3(x+2)(x-1) & = & 0\\
x=-2 & \mbox{ and } & x=1
\end{eqnarray*}
\newpage
\end{enumerate}

\subsubsection*{Solve the equation using the quadratic formula. ex. $3x^{2}=-2x+2$.}
\begin{enumerate}
\item Set the equation equal to $0$.
\[
3x^{2}+2x-2=0
\]
\item Identify the $a,b,c$ in $ax^{2}+bx+c=0$
\[
a=3,\,b=2,\,c=-2
\]
\item Use the quadratic formula, $x=\frac{-b\pm\sqrt{b^{2}-4ac}}{2a}$,
and simplify.
\[
x=\frac{-2\pm\sqrt{(2)^{2}-4(3)(-2)}}{2(3)}=\frac{-2\pm\sqrt{28}}{6}=\frac{-1\pm\sqrt{7}}{3}
\]
\end{enumerate}

\subsubsection*{Solve the equation by completing the square. ex. $x^{2}=-4x+2$.}
\begin{enumerate}
\item Set the equation eqaul to ``$c$'', in this case $2$, and divide
by the ``$a$''.
\begin{eqnarray*}
x^{2}+4x & = & 2\\
x^{2}+4x & = & 2
\end{eqnarray*}
\item Add $\left(\frac{b}{2}\right)^{2}$ to both sides; in our example
$b=4$ so add $\left(\frac{4}{2}\right)^{2}=\left(2\right)^{2}=4$.
\[
x^{2}+4x+4=2+4
\]
\item Factor, and solve (note it will always be $(x+\frac{b}{2})^{2}$)
\begin{eqnarray*}
(x+2)^{2} & = & 6\\
x+2 & =\pm & \sqrt{6}\\
x & =-2\pm & \sqrt{6}
\end{eqnarray*}
\end{enumerate}

\subsubsection*{Use the 5 step strategy for solving the word problem. When 3 times
the number is subtracted 7 times the number, the result is 44.}
\begin{enumerate}
\item Summarize: ``3 times a number'' means, $3x$. ``7 times a number'',
means $7x$. ``Subtracted from'' means, $7x-3x$. ``Result is 44''
means, $=44$.
\item Write an equation: $7x-3x=44$
\item Solve: $x=11$.
\item Check: $7(11)-3(11)=44\checkmark$
\item State your answer as a sentence (this isn't applicable here).
\end{enumerate}
\textbf{Solve the word problem. ex. }A car rental place charges \$300
a week to rent a car plus an additional \$0.50 per mile. How many
miles can you drive in a week for \$400.
\begin{enumerate}
\item Write as an equation.
\[
300+0.50x=400
\]
\item Solve. 
\begin{eqnarray*}
0.50x & = & 100\\
x & = & \frac{100}{0.50}=200\mbox{ miles}
\end{eqnarray*}
\end{enumerate}

\end{document}
